\documentclass[11pt]{article}

\usepackage[margin=1in]{geometry}
\usepackage[T1]{fontenc}
\usepackage{amsmath,amsthm,mathtools}
\usepackage{newtxtext,newtxmath}
\usepackage{microtype}
\usepackage{setspace}
\usepackage{booktabs,array,tabularx}
\usepackage{enumitem}
\usepackage{tikz}
\usepackage{natbib}
\usepackage[hidelinks]{hyperref}
\hypersetup{pdftitle={Developer Responses to New York City's 100-Unit Threshold: A Notch-and-Kink Model},pdfauthor={Jacob Herbstman}}
\usepackage{fancyhdr}
\usepackage{xcolor}

\setstretch{1.06}
\setlength{\parindent}{1.25em}
\setlength{\parskip}{0.15em}
\setlist[itemize]{leftmargin=1.4em,itemsep=0.15em,topsep=0.25em}
\setlist[enumerate]{leftmargin=1.6em,itemsep=0.15em,topsep=0.25em}

\newtheorem{proposition}{Proposition}
\newtheorem{assumption}{Assumption}
\newtheorem{remark}{Remark}

\newcommand{\ind}{\mathbf{1}}
\newcommand{\E}{\mathbb{E}}
\newcommand{\R}{\mathbb{R}}
\newcommand{\dd}{\,\mathrm{d}}

\pagestyle{fancy}
\fancyhf{}
\fancyhead[L]{\small Developer Responses to NYC's 100-Unit Threshold}
\fancyhead[R]{\small \thepage}
\renewcommand{\headrulewidth}{0.3pt}

\title{\vspace{-1.25cm}\textbf{Developer Responses to New York City's 100-Unit Threshold}\\[0.35em]
\large A Notch-and-Kink Model}
\author{Jacob Herbstman}
\date{August 2026}

\begin{document}
\maketitle
\vspace{-1.2em}

\begin{abstract}
This note develops a parsimonious model of residential developer behavior around New York City's 100-unit threshold under the 485-x tax incentive. Crossing the threshold may create both a discontinuous project-level burden---a notch---and a persistent increase in marginal project cost---a kink. The model predicts a mass point at 99 units, a missing source interval in the no-policy unit-count distribution, and horizontal compression of projects that remain above the threshold. A reweighted historical distribution identifies the no-policy size of the marginal bunching project from excess mass at 99; the surviving upper tail identifies, or bounds, the compression parameter. Given a curvature parameter, these behavioral objects map into a marginal-cost wedge, a normalized level notch, and the total policy burden faced by the marginal developer. The model also yields transparent implications for housing units, total floor area, and average unit size.
\end{abstract}

\section{Motivation}

New York City's 485-x program changes the regulatory treatment of rental projects at 100 units. Relative to projects below 100 units, large rental projects face different affordability requirements and construction compensation rules; projects at 150 units in designated areas face an additional regime change \citep{nys_485x,nyc_hpd_485x}. The sharp empirical motivation is a large increase in projects proposed at exactly 99 units after the policy took effect.

The model combines two ideas. First, size-dependent regulations can create a discontinuity in total project value or cost, generating bunching below a threshold and missing mass above it \citep{klevenwaseem2013,bestkleven2018,kopczukmunroe2015}. Second, the high-regime cost schedule may remain persistently steeper, compressing the surviving upper tail rather than allowing it to return point-by-point to its no-policy distribution \citep{gourioroys2014,garicano2016}. The setup follows the isoelastic cost structure used by \citet{brueckner2024}, but the policy combines a notch and a kink rather than a pure kink.

\section{Developer problem}

\subsection{No-policy benchmark}

A developer chooses the number of residential units, $n$. Unit counts are discrete in the data; calculus is used to characterize choices within each regime. Let $v>0$ denote a project-specific marginal profitability shifter, net of ordinary linear per-unit costs. Smooth no-policy profit is
\begin{equation}
\pi_0(n;v)=vn-C_0(n),
\qquad
C_0(n)=\frac{\alpha}{1+\lambda}n^{1+\lambda},
\label{eq:baseline_profit}
\end{equation}
where $\alpha>0$ and $\lambda>0$. Marginal cost is $C_0'(n)=\alpha n^\lambda$, so the unique no-policy optimum is
\begin{equation}
 n_0(v)=\left(\frac{v}{\alpha}\right)^{1/\lambda}.
\label{eq:n0}
\end{equation}
The parameter $1/\lambda$ is the elasticity of desired unit count with respect to $v$.

\subsection{Policy cost schedule}

Let $N=100$ be the first unit count in the high regime and $b=N-1=99$ the largest avoiding choice. For $n\geq N$, define the incremental net policy burden
\begin{equation}
D(n)=K+\frac{\beta}{1+\lambda}\left(n^{1+\lambda}-N^{1+\lambda}\right),
\qquad K\geq0,\quad \beta\geq0.
\label{eq:policy_burden}
\end{equation}
Total profit is
\begin{equation}
\pi(n;v)=\pi_0(n;v)-\ind\{n\geq N\}D(n).
\label{eq:full_profit}
\end{equation}
The parameter $K=D(N)$ is the net level discontinuity at entry. It can collect threshold-triggered wage, affordability, compliance, and tax-benefit effects. The parameter $\beta$ raises marginal cost above the threshold because
\begin{equation}
C_H'(n)=(\alpha+\beta)n^\lambda,
\qquad n>N.
\end{equation}
Thus $\beta=0$ gives a pure notch, $K=0$ gives the pure-kink structure in \citet{brueckner2024}, and $K,\beta>0$ gives a notch plus a kink.

\subsection{Choice conditional on entering}

Conditional on entering the high regime, the developer solves the constrained problem $\max_{n\geq N}\pi_H(n;v)$. The unconstrained high-regime first-order condition is
\begin{equation}
 v=(\alpha+\beta)n^\lambda.
\end{equation}
Define the compression factor
\begin{equation}
\theta\equiv\left(\frac{\alpha+\beta}{\alpha}\right)^{1/\lambda}\geq1.
\label{eq:theta}
\end{equation}
Combining the low- and high-regime first-order conditions gives the unconstrained high-regime optimum $n_0(v)/\theta$. Respecting the lower boundary of the high regime, the best feasible high-regime choice is therefore
\begin{equation}
 n_H(v)=\max\left\{N,\frac{n_0(v)}{\theta}\right\}.
\label{eq:compression}
\end{equation}
Thus $\theta=1$ under a pure notch, while $\theta>1$ horizontally compresses projects that still enter. Equation \eqref{eq:compression} only characterizes the best choice conditional on entry; a developer whose best high-regime option is $N$ may still prefer $b=99$ once the level discontinuity $K$ is taken into account.

\subsection{Entry and the marginal developer}

For a developer with $n_0(v)>b$, the best avoiding choice is $b$. Let
\begin{equation}
 V_L(v)=vb-C_0(b),
 \qquad
 V_H(v)=\max_{n\geq N}\pi_H(n;v),
\end{equation}
and define the gain from entry as $\Delta(v)=V_H(v)-V_L(v)$. The envelope theorem gives the monotone-sorting result directly. Since $V_H(v)=\pi_H(n_H(v);v)$,
\begin{align}
 V_H'(v)
 &=\frac{\partial\pi_H}{\partial v}
 +\frac{\partial\pi_H}{\partial n}n_H'(v)
 =n_H(v),
 \label{eq:envelope_main}
\end{align}
where the second term vanishes by the first-order condition at an interior optimum and because $n_H(v)=N$ is locally constant when the boundary binds. Since $V_L'(v)=b$,
\begin{equation}
\boxed{\Delta'(v)=n_H(v)-b>0.}
\label{eq:monotone_entry}
\end{equation}
There is therefore a unique marginal type $v^\dagger$ satisfying $\Delta(v^\dagger)=0$. Define
\begin{equation}
 q\equiv n_0(v^\dagger),
 \qquad
 h\equiv n_H(v^\dagger)
 =\max\left\{N,\frac{q}{\theta}\right\}.
\label{eq:q_h}
\end{equation}
The resulting decision rule is
\begin{equation}
 n^*(v)=
 \begin{cases}
 n_0(v), & n_0(v)\leq b,\\[3pt]
 b, & b<n_0(v)<q,\\[3pt]
 \max\left\{N,n_0(v)/\theta\right\}, & n_0(v)\geq q.
 \end{cases}
\label{eq:decision_rule}
\end{equation}
The model therefore permits a boundary choice at $N$. If $q<\theta N$, a positive interval of developers enters but chooses exactly $N$. If instead $q\geq\theta N$, the boundary is weakly dominated by $b=99$ for every type whose conditional high-regime optimum is $N$; under strict inequality, the first crosser is interior and the surviving upper tail begins at $q/\theta>N$. The current data contain no mass at exactly 100 and little mass immediately above it, so the main analysis emphasizes this strict interior-crossing branch, while retaining boundary entry as a technically feasible model outcome. Figure \ref{fig:mapping} illustrates the interior-crossing branch.

\begin{figure}[t]
\centering
\begin{tikzpicture}[x=1.05cm,y=0.78cm,>=stealth]
  % axes
  \draw[->] (0,0) -- (8.2,0) node[right] {$x$: no-policy units};
  \draw[->] (0,0) -- (0,6.2) node[above] {$n^*(x)$};
  % 45 degree lower branch
  \draw[thick] (0.45,0.45) -- (2.15,2.15);
  % bunching branch
  \draw[thick] (2.85,2.15) -- (5.30,2.15);
  % upper compressed branch
  \draw[thick] (5.30,3.55) -- (7.65,5.12);
  % dashed guides
  \draw[dashed] (2.15,0) -- (2.15,2.15);
  \draw[dashed] (2.85,0) -- (2.85,2.15);
  \draw[dashed] (5.30,0) -- (5.30,3.55);
  \draw[dashed] (0,2.15) -- (5.30,2.15);
  \draw[dashed] (0,3.55) -- (5.30,3.55);
  % labels
  \node[below] at (2.15,0) {$b$};
  \node[below] at (2.85,0) {$N$};
  \node[below] at (5.30,0) {$q$};
  \node[left] at (0,2.15) {$b$};
  \node[left] at (0,3.55) {$q/\theta$};
  \node[align=center] at (4.05,1.45) {move to\\99};
  \node[align=center,rotate=33] at (6.70,4.72) {slope $1/\theta$};
  \node[align=center] at (6.15,2.65) {gap};
\end{tikzpicture}
\caption{Schematic policy mapping in the interior-crossing branch $q>\theta N$. Projects with no-policy size between $N$ and $q$ move to $b=99$; projects above $q$ remain in the high regime but are compressed by $1/\theta$. The general model also permits a boundary choice at $N$.}
\label{fig:mapping}
\end{figure}

\section{Distributional implications and identification}

Let $H_0(j)$ denote the expected number of projects choosing exactly $j$ units without the policy, and let $h_0(x)$ denote a continuous interpolation of this distribution. The model generates two distinct moments.

\subsection{Bunching identifies the no-policy source frontier}

Excess mass at 99 is
\begin{equation}
 B\equiv H_{\mathrm{obs}}(b)-H_0(b).
\end{equation}
Under \eqref{eq:decision_rule}, these projects come from no-policy choices between $N$ and $q$. The exact integer and continuous equations are
\begin{equation}
 B=\sum_{j=N}^{q-1}H_0(j)
 \qquad\text{and}\qquad
 B=\int_N^q h_0(x)\dd x.
\label{eq:bunching_mass}
\end{equation}
Given a counterfactual $H_0$, the spike identifies $q$ by accumulating mass from 100 upward. If the frontier cuts through an integer bin, the empirical implementation allocates that terminal bin fractionally. Importantly, $\theta$ does not enter \eqref{eq:bunching_mass}: the spike identifies which no-policy projects leave the high regime, not how the surviving projects are compressed.

\subsection{The upper tail identifies the compression factor}

The model permits two high-regime components. If $q<\theta N$, developers with no-policy sizes between $q$ and $\theta N$ choose the boundary $N$, generating an atom
\begin{equation}
 B_N=\int_q^{\theta N}h_0(x)\dd x.
\label{eq:boundary_atom}
\end{equation}
This atom is zero when $q\geq\theta N$. For the continuous part of the upper tail, observed size is $n=x/\theta$, where $x$ is no-policy size. Conservation of mass implies
\begin{equation}
 h_1(n)=\theta h_0(\theta n),
 \qquad n>\max\left\{N,\frac{q}{\theta}\right\}.
\label{eq:tail_density}
\end{equation}
Regardless of whether the boundary branch is active, observed projects between $N$ and any $m\geq h$ originate from no-policy projects between $q$ and $\theta m$:
\begin{equation}
 H_{\mathrm{obs}}([N,m])=\int_q^{\theta m}h_0(x)\dd x.
\label{eq:cumulative_theta}
\end{equation}
Setting $m=q$ yields a particularly transparent identifying moment:
\begin{equation}
 H_{\mathrm{obs}}([N,q])=\int_q^{\theta q}h_0(x)\dd x.
\label{eq:theta_moment}
\end{equation}
Thus $q$ is identified by the mass moved to 99, while $\theta$ is identified, in principle, by the mass and shape of the surviving upper tail, including any boundary atom. In sparse samples, \eqref{eq:theta_moment} can be used for a point estimate and the full distributional restrictions \eqref{eq:boundary_atom}--\eqref{eq:tail_density} for minimum-distance or likelihood estimation. A sensitivity analysis over $\theta$ remains useful.

The condition
\begin{equation}
 q>\theta N
 \quad\Longleftrightarrow\quad
 \theta<\frac{q}{N}
\label{eq:theta_bound}
\end{equation}
characterizes the interior-crossing branch; it is not imposed by the model. If instead $q<\theta N$, the model predicts positive mass at exactly $N$. The absence of observed 100-unit projects therefore provides direct evidence about which branch is empirically relevant.

\subsection{Why the missing-mass curve need not recover}

Define the cumulative deficit through $m$ as
\begin{equation}
 \mathcal{D}(m)=\int_N^m\left[h_0(n)-h_1(n)\right]\dd n.
\end{equation}
For $m\geq q$, equations \eqref{eq:bunching_mass} and \eqref{eq:tail_density} imply
\begin{equation}
 \boxed{\mathcal{D}(m)=B-\int_m^{\theta m}h_0(x)\dd x.}
\label{eq:deficit_identity}
\end{equation}
The second term is no-policy mass from above $m$ compressed into the observed range below $m$. Under a pure notch, $\theta=1$ and the term is zero, so the deficit reaches $B$. With a kink, $\theta>1$ and pointwise recovery is not expected.

\section{Mapping behavior into policy costs}

The model distinguishes three related cost objects.

\subsection{Marginal-cost wedge}

From \eqref{eq:theta},
\begin{equation}
 \boxed{\frac{\beta}{\alpha}=\theta^\lambda-1.}
\label{eq:beta_wedge}
\end{equation}
This is the proportional increase in marginal cost at a fixed project size above the threshold.

\subsection{Level notch}

The marginal developer is indifferent between the avoiding choice $b$ and the optimal high-regime choice $h=\max\{N,q/\theta\}$. Hence
\begin{equation}
\begin{aligned}
K
={}&v^\dagger(h-b)-\left[C_0(h)-C_0(b)\right]\\
&-\frac{\beta}{1+\lambda}\left(h^{1+\lambda}-N^{1+\lambda}\right).
\end{aligned}
\label{eq:K_main}
\end{equation}
The first term is the gross value of expanding beyond 99, the second is ordinary incremental cost, and the third is the accumulated kink cost above 100. Use
\begin{equation}
 v^\dagger=\alpha q^\lambda,
 \qquad
 \frac{\beta}{\alpha}=\theta^\lambda-1,
\end{equation}
and define the dimensionless ratios
\begin{equation}
 x\equiv\frac{q}{b},
 \qquad
 r\equiv\frac{N}{b},
 \qquad
 y\equiv\frac{h}{b}=\max\left\{r,\frac{x}{\theta}\right\},
 \qquad
 \kappa\equiv\frac{K}{C_0(b)}.
\end{equation}
Dividing \eqref{eq:K_main} by $C_0(b)=\alpha b^{1+\lambda}/(1+\lambda)$ removes the unidentified cost scale $\alpha$ and yields the general normalized notch:
\begin{equation}
\boxed{
\kappa
=(1+\lambda)x^\lambda(y-1)
-\left(y^{1+\lambda}-1\right)
-(\theta^\lambda-1)\left(y^{1+\lambda}-r^{1+\lambda}\right).}
\label{eq:kappa_general}
\end{equation}
In the interior-crossing branch, $y=x/\theta$, and \eqref{eq:kappa_general} simplifies to
\begin{equation}
 \boxed{
 \kappa
 =1+\frac{\lambda}{\theta}x^{1+\lambda}
 -(1+\lambda)x^\lambda
 +(\theta^\lambda-1)r^{1+\lambda}.}
\label{eq:kappa}
\end{equation}
If the marginal crosser instead chooses the boundary, $y=r$ and
\begin{equation}
 \kappa
 =(1+\lambda)x^\lambda(r-1)-\left(r^{1+\lambda}-1\right).
\label{eq:kappa_boundary}
\end{equation}
The pure-notch special case $\theta=1$ is
\begin{equation}
 \kappa_{\mathrm{notch}}
 =1+\lambda x^{1+\lambda}-(1+\lambda)x^\lambda.
\label{eq:pure_notch_kappa}
\end{equation}
For fixed $(q,\lambda)$, $\kappa$ falls with $\theta$ within the interior-crossing branch: a stronger marginal-cost effect requires a smaller level jump to generate the same bunching frontier.

\subsection{Total burden faced by the marginal crosser}

The total policy burden at $h$ is
\begin{equation}
D(h)=K+\frac{\beta}{1+\lambda}\left(h^{1+\lambda}-N^{1+\lambda}\right).
\end{equation}
Normalized by $C_0(b)$, it can be written generally as
\begin{equation}
 \boxed{
 \delta\equiv\frac{D(h)}{C_0(b)}
 =(1+\lambda)x^\lambda(y-1)-\left(y^{1+\lambda}-1\right),}
\label{eq:delta}
\end{equation}
where $y=\max\{r,x/\theta\}$. Equivalently,
\begin{equation}
 \delta=\kappa+(\theta^\lambda-1)\left(y^{1+\lambda}-r^{1+\lambda}\right).
\end{equation}
The three objects answer different questions:
\begin{center}
\begin{tabularx}{0.96\textwidth}{@{}>{\raggedright\arraybackslash}p{0.22\textwidth} >{\centering\arraybackslash}p{0.25\textwidth} X@{}}
\toprule
Object & Formula & Interpretation \\
\midrule
Marginal-cost wedge & $\beta/\alpha=\theta^\lambda-1$ & Percentage increase in marginal cost in the high regime. \\
Level notch & $\kappa=K/C_0(b)$ & Net discontinuity triggered at 100, relative to modeled smooth cost at 99. \\
Total marginal-project burden & $\delta=D(h)/C_0(b)$ & Total direct policy burden paid by the marginal crosser at its optimal high-regime size. \\
\bottomrule
\end{tabularx}
\end{center}
The denominator $C_0(b)$ is the explicitly modeled convex cost component. Converting these objects into a share of all-in project cost, or into dollars, requires additional cost information.

\section{Empirical implementation}

The baseline counterfactual need not assign a latent no-policy unit count to each post-policy lot. Let historical parent projects retain their observed untreated unit counts $N_j^0$ and receive calibration weights $w_j$ chosen using predetermined site characteristics. The transported no-policy distribution is
\begin{equation}
 \widehat H_0(n)=\sum_{j\in\mathrm{pre}}w_j\ind\{N_j^0=n\}.
\label{eq:transport}
\end{equation}
This preserves historical exact-count heaping and estimates an aggregate distribution for a population with the post-policy site composition.

A natural estimation sequence is:
\begin{enumerate}
\item Estimate natural mass at 99 and excess bunching $\widehat B$.
\item Solve \eqref{eq:bunching_mass} for the no-policy frontier $\widehat q$.
\item Estimate $\theta$ from \eqref{eq:theta_moment} or the full transformed distribution \eqref{eq:boundary_atom}--\eqref{eq:tail_density}. The general estimator should allow the boundary branch, while a sensitivity grid can emphasize values below $\widehat q/N$ when the data contain no mass at 100.
\item For each $(\theta,\lambda)$, report $\beta/\alpha$, $\kappa$, and $\delta$ using \eqref{eq:beta_wedge}, \eqref{eq:kappa_general}, and \eqref{eq:delta}.
\end{enumerate}
The 150-unit threshold creates a separate policy regime. Empirical moments used to estimate the 100-unit model should therefore remain below 150 or explicitly model the second threshold. When $\widehat q$ is near 122, values $\theta<1.22$ place the marginal crosser above 100 and keep the identifying interval $[q,\theta q]$ below 150. Values at or above this boundary predict an atom at 100 and are directly testable against the observed histogram.

The parameter $\lambda$ governs curvature in unit-count space. It should not be imported mechanically from a floor-space production function if developers can respond by changing apartment size rather than total floor area. The preferred approach is to report cost mappings over a plausible range of $\lambda$ until project cost or pro forma data provide additional discipline.

\section{Quantity and design implications}

The model separates unit losses among bunchers from changes among crossers. Aggregate units suppressed among bunchers are
\begin{equation}
 \Delta U_B=\int_N^q(x-b)h_0(x)\dd x,
\end{equation}
while losses among crossers are, in the general model,
\begin{equation}
 \Delta U_H=\int_q^\infty\left[x-\max\left\{N,\frac{x}{\theta}\right\}\right]h_0(x)\dd x.
\end{equation}
In the interior-crossing branch this simplifies to $\int_q^\infty x(1-1/\theta)h_0(x)\dd x$. The buncher term is identified from the local counterfactual below $q$; the crosser term requires the broader upper tail and, in this application, an extension through the 150-unit regime.

Let total residential floor area satisfy the reduced-form relationship $Q^*(n,z)=A(z)n^\eta$. Average floor area per unit is $s=Q/n$. A project that bunches from no-policy size $x$ to $b$ satisfies
\begin{equation}
 \frac{Q_1}{Q_0}=\left(\frac{b}{x}\right)^\eta,
 \qquad
 \frac{s_1}{s_0}=\left(\frac{x}{b}\right)^{1-\eta}.
\label{eq:floor_area_buncher}
\end{equation}
For a crosser with observed size $n_H(x)=\max\{N,x/\theta\}$,
\begin{equation}
 \frac{Q_1}{Q_0}=\left(\frac{n_H(x)}{x}\right)^\eta,
 \qquad
 \frac{s_1}{s_0}=\left(\frac{x}{n_H(x)}\right)^{1-\eta}.
\label{eq:floor_area_crosser}
\end{equation}
In the interior branch these ratios simplify to $\theta^{-\eta}$ and $\theta^{1-\eta}$.
If $\eta<1$, unit counts fall more than floor area, so average apartments become larger. This provides auxiliary predictions for floor area, FAR, square feet per unit, and bedroom mix.

\section{Discussion}

The baseline is intentionally parsimonious. It assumes a common $(K,\beta)$, a fixed set of development opportunities, and monotone sorting into bunching and crossing. Boundary entry at 100 is technically permitted and generates a testable atom; the current distribution strongly favors the interior-crossing branch but this is an empirical implication rather than a primitive restriction. The main extensions are heterogeneous policy burdens, project delay or cancellation, participation outside 485-x, alternative avoidance points such as 95 units, and an explicit model of the 150-unit threshold. These extensions should be introduced only when the baseline distributional restrictions fail in ways that are both economically and statistically meaningful.

\appendix
\section{Derivations}

\subsection{Monotone entry}

For types with $n_0(v)>b$, the low-regime value is $V_L(v)=vb-C_0(b)$, so $V_L'(v)=b$. Optimized high-regime value is $V_H(v)=\pi_H(n_H(v);v)$. The multivariable chain rule gives
\begin{equation}
 V_H'(v)=\frac{\partial\pi_H}{\partial v}
 +\frac{\partial\pi_H}{\partial n}n_H'(v).
\end{equation}
At an interior optimum, the second term is zero by the first-order condition. When the boundary binds, $n_H(v)=N$ is locally constant, so $n_H'(v)=0$. In either case, $\partial\pi_H/\partial v=n_H(v)$ and $V_H'(v)=n_H(v)$ almost everywhere. Hence
\begin{equation}
\Delta'(v)=V_H'(v)-V_L'(v)=n_H(v)-b>0.
\end{equation}

\subsection{Upper-tail density and cumulative deficit}

For interior crossers, $n=x/\theta$, or $x=\theta n$. A small observed interval $\dd n$ corresponds to no-policy width $\dd x=\theta\dd n$. Conservation of mass gives
\begin{equation}
 h_1(n)\dd n=h_0(\theta n)\theta\dd n,
\end{equation}
which yields the continuous-tail restriction \eqref{eq:tail_density}. If $q<\theta N$, the no-policy interval $[q,\theta N]$ maps to the boundary $N$, producing \eqref{eq:boundary_atom}.

For any $m\geq q$, the total observed high-regime mass through $m$ is $\int_q^{\theta m}h_0(x)\dd x$, including any boundary atom. Hence
\begin{align}
\mathcal{D}(m)
&=\int_N^m h_0(n)\dd n-\int_q^{\theta m}h_0(x)\dd x\\
&=\int_N^q h_0(x)\dd x-\int_m^{\theta m}h_0(x)\dd x,
\end{align}
which is \eqref{eq:deficit_identity}.

\subsection{Normalized level-notch formula}

The marginal developer's indifference condition is \eqref{eq:K_main}. Use
\begin{equation}
 v^\dagger=\alpha q^\lambda,
 \qquad
 h=yb,
 \qquad
 y=\max\left\{r,\frac{x}{\theta}\right\},
 \qquad
 \frac{\beta}{\alpha}=\theta^\lambda-1.
\end{equation}
Divide by $C_0(b)=\alpha b^{1+\lambda}/(1+\lambda)$. The three normalized terms are
\begin{align}
\frac{v^\dagger(h-b)}{C_0(b)}
&=(1+\lambda)x^\lambda(y-1),\\
\frac{C_0(h)-C_0(b)}{C_0(b)}
&=y^{1+\lambda}-1,\\
\frac{\beta(h^{1+\lambda}-N^{1+\lambda})/(1+\lambda)}{C_0(b)}
&=(\theta^\lambda-1)\left(y^{1+\lambda}-r^{1+\lambda}\right).
\end{align}
Subtracting the latter two terms from the first yields \eqref{eq:kappa_general}. In the interior branch $y=x/\theta$, and collecting powers of $x$ and $\theta$ yields \eqref{eq:kappa}.

\subsection{Special cases}

Under a pure notch, $\beta=0$ and $\theta=1$, producing \eqref{eq:pure_notch_kappa}. If $\lambda=1$, then
\begin{equation}
\kappa_{\mathrm{notch}}=\left(\frac{q-b}{b}\right)^2.
\end{equation}

If the high regime proportionally raises the entire smooth cost function by $\tau$, then $C_H(n)=(1+\tau)C_0(n)$. This is nested by setting
\begin{equation}
\beta=\tau\alpha,
\qquad
K=\tau C_0(N).
\end{equation}
The policy burden then satisfies $D(n)=\tau C_0(n)$ for all $n\geq N$.

\begin{thebibliography}{99}

\bibitem[Best and Kleven(2018)]{bestkleven2018}
Best, Michael Carlos, and Henrik Jacobsen Kleven. 2018.
``Housing Market Responses to Transaction Taxes: Evidence from Notches and Stimulus in the UK.''
\emph{Review of Economic Studies} 85(1): 157--193.

\bibitem[Brueckner et al.(2024)]{brueckner2024}
Brueckner, Jan K., David Leather, and Miguel Zerecero. 2024.
``Bunching in Real-Estate Markets: Regulated Building Heights in New York City.''
\emph{Journal of Urban Economics} 143: 103683.

\bibitem[Garicano et al.(2016)]{garicano2016}
Garicano, Luis, Claire Lelarge, and John Van Reenen. 2016.
``Firm Size Distortions and the Productivity Distribution: Evidence from France.''
\emph{American Economic Review} 106(11): 3439--3479.

\bibitem[Gourio and Roys(2014)]{gourioroys2014}
Gourio, Fran\c{c}ois, and Nicolas Roys. 2014.
``Size-Dependent Regulations, Firm Size Distribution, and Reallocation.''
\emph{Quantitative Economics} 5(2): 377--416.

\bibitem[Kleven and Waseem(2013)]{klevenwaseem2013}
Kleven, Henrik J., and Mazhar Waseem. 2013.
``Using Notches to Uncover Optimization Frictions and Structural Elasticities: Theory and Evidence from Pakistan.''
\emph{Quarterly Journal of Economics} 128(2): 669--723.

\bibitem[Kopczuk and Munroe(2015)]{kopczukmunroe2015}
Kopczuk, Wojciech, and David Munroe. 2015.
``Mansion Tax: The Effect of Transfer Taxes on the Residential Real Estate Market.''
\emph{American Economic Journal: Economic Policy} 7(2): 214--257.

\bibitem[New York City HPD(2026)]{nyc_hpd_485x}
New York City Department of Housing Preservation and Development. 2026.
``485-x: Affordable Neighborhoods for New Yorkers.''
\url{https://www.nyc.gov/site/hpd/services-and-information/tax-incentives-485-x.page}.
Accessed August 2026.

\bibitem[New York State Legislature(2024)]{nys_485x}
New York State Legislature. 2024.
``S. 8306-C / A. 8806-C, Part R: Affordable Neighborhoods for New Yorkers Tax Incentive.''
\url{https://legislation.nysenate.gov/pdf/bills/2023/s8306c}.

\end{thebibliography}

\end{document}
